% GENERATED FILE. Edit canonical lessons, then run npm run generate:books.
% canonical-source-hash: fnv1a64:41e09de9beeabe6f
% canonical-lessons: FR-C294-interpeller, FR-C294-arrestation, FR-C294-menottes, FR-C294-garde-a-vue, FR-C294-interroger

\chapter{L'arrestation}
\label{ch:fr-l-arrestation}
\hlchaptermodality{294}

\begin{chapteropening}
\textbf{By the end of this chapter:} \emph{I can describe an arrest: stopping someone, the arrest, handcuffs, police custody and questioning.}

Complete the last lesson of chapter 294: I can describe an arrest: stopping someone, the arrest, handcuffs, police custody and questioning.
\end{chapteropening}

\section[interpeller]{interpeller — to stop and question (police); to call out to}

\label{lesson:FR-C294-interpeller}



\begin{quote}
Before the new one: say the French for the alibi, then the French for the suspicion.
\end{quote}

\subsection*{You'll want to know: interpeller}
\textbf{interpeller} — \textquotedblleft{}to stop and question (police); to call out to\textquotedblright{}.

\begin{quote}
La police a interpellé deux jeunes hommes.
\end{quote}

\begin{quote}
\emph{The police stopped two young men for questioning.}
\end{quote}

In the news, \emph{interpeller} means the police took someone in. In the street, it means to call out to someone.

\begin{cousinweb}
\emph{interpeller} is a learned borrowing from Latin \emph{interpellāre} \textquotedblleft{}to interrupt, to address\textquotedblright{} (EtymDB).
\end{cousinweb}

\begin{grammarlens}[title={What you've built}]
One more word for telling, explaining and describing at B1.
\end{grammarlens}

\subsection*{Your turn}
\begin{itemize}
\raggedright
  \item \emph{Say it:} \emph{interpeller}
  \item \emph{Say it:} \emph{interpeller}, once more
  \item \emph{Say it:} a sentence of your own with \emph{interpeller}
  \item \emph{Recall:} say the French for the alibi, then the French for the suspicion, then say \emph{interpeller} again
\end{itemize}

\begin{tcolorbox}[breakable,colback=teal!4,colframe=teal!35!black,title={Before you move on}]
What does interpeller mean? (\textquotedblleft{}To stop and question (police); to call out to\textquotedblright{}.) Say it once more.
\end{tcolorbox}
\section[l'arrestation]{l'arrestation — the arrest}

\label{lesson:FR-C294-arrestation}



\begin{quote}
Before the new one: say the French for the suspicion, then the French for to stop and question.
\end{quote}

\subsection*{You'll want to know: l'arrestation}
\textbf{l'arrestation} — \textquotedblleft{}the arrest\textquotedblright{}.

\begin{quote}
L'arrestation a eu lieu tôt le matin.
\end{quote}

\begin{quote}
\emph{The arrest took place early in the morning.}
\end{quote}

The verb is \emph{arrêter}: \emph{il a été arrêté}.

\begin{cousinweb}
\emph{arrestation} comes from Medieval Latin \emph{arrestātiō} (EtymDB), built on \emph{arrestāre} \textquotedblleft{}to stop\textquotedblright{}, from \emph{ad} and \emph{restāre} \textquotedblleft{}to stay back\textquotedblright{} (Lewis \& Short).
\end{cousinweb}

\begin{grammarlens}[title={What you've built}]
One more word for telling, explaining and describing at B1.
\end{grammarlens}

\subsection*{Your turn}
\begin{itemize}
\raggedright
  \item \emph{Say it:} \emph{l'arrestation}
  \item \emph{Say it:} \emph{l'arrestation}, once more
  \item \emph{Say it:} a sentence of your own with \emph{l'arrestation}
  \item \emph{Recall:} say the French for the suspicion, then the French for to stop and question, then say \emph{l'arrestation} again
\end{itemize}

\begin{tcolorbox}[breakable,colback=teal!4,colframe=teal!35!black,title={Before you move on}]
What does l'arrestation mean? (\textquotedblleft{}The arrest\textquotedblright{}.) Say it once more.
\end{tcolorbox}
\section[les menottes]{les menottes — the handcuffs}

\label{lesson:FR-C294-menottes}



\begin{quote}
Before the new one: say the French for to stop and question, then the French for the arrest.
\end{quote}

\subsection*{You'll want to know: les menottes}
\textbf{les menottes} — \textquotedblleft{}the handcuffs\textquotedblright{}.

\begin{quote}
Le policier lui a mis les menottes.
\end{quote}

\begin{quote}
\emph{The police officer put the handcuffs on him.}
\end{quote}

Literally 'little hands', from \emph{main}.

\begin{cousinweb}
\emph{menottes} is a diminutive of \emph{main} \textquotedblleft{}hand\textquotedblright{}, which is inherited from Latin \emph{manus} \textquotedblleft{}hand\textquotedblright{} (Lewis \& Short).
\end{cousinweb}

\begin{grammarlens}[title={What you've built}]
One more word for telling, explaining and describing at B1.
\end{grammarlens}

\subsection*{Your turn}
\begin{itemize}
\raggedright
  \item \emph{Say it:} \emph{les menottes}
  \item \emph{Say it:} \emph{les menottes}, once more
  \item \emph{Say it:} a sentence of your own with \emph{les menottes}
  \item \emph{Recall:} say the French for to stop and question, then the French for the arrest, then say \emph{les menottes} again
\end{itemize}

\begin{tcolorbox}[breakable,colback=teal!4,colframe=teal!35!black,title={Before you move on}]
What does les menottes mean? (\textquotedblleft{}The handcuffs\textquotedblright{}.) Say it once more.
\end{tcolorbox}
\section[la garde à vue]{la garde à vue — police custody}

\label{lesson:FR-C294-garde-a-vue}



\begin{quote}
Before the new one: say the French for the arrest, then the French for the handcuffs.
\end{quote}

\subsection*{You'll want to know: la garde à vue}
\textbf{la garde à vue} — \textquotedblleft{}police custody\textquotedblright{}.

\begin{quote}
Le jeune homme a passé la nuit en garde à vue.
\end{quote}

\begin{quote}
\emph{The young man spent the night in police custody.}
\end{quote}

\emph{Être placé en garde à vue} — to be held for questioning.

\begin{cousinweb}
\emph{garde} comes from \emph{garder} \textquotedblleft{}to keep, to watch\textquotedblright{}, from a Frankish verb \emph{wardōn} (EtymDB); \emph{vue} is built on \emph{voir}, from Latin \emph{vidēre} \textquotedblleft{}to see\textquotedblright{}: kept in sight.
\end{cousinweb}

\begin{grammarlens}[title={What you've built}]
One more word for telling, explaining and describing at B1.
\end{grammarlens}

\subsection*{Your turn}
\begin{itemize}
\raggedright
  \item \emph{Say it:} \emph{la garde à vue}
  \item \emph{Say it:} \emph{la garde à vue}, once more
  \item \emph{Say it:} a sentence of your own with \emph{la garde à vue}
  \item \emph{Recall:} say the French for the arrest, then the French for the handcuffs, then say \emph{la garde à vue} again
\end{itemize}

\begin{tcolorbox}[breakable,colback=teal!4,colframe=teal!35!black,title={Before you move on}]
What does la garde à vue mean? (\textquotedblleft{}Police custody\textquotedblright{}.) Say it once more.
\end{tcolorbox}
\section[interroger]{interroger — to question, to interrogate}

\label{lesson:FR-C294-interroger}



\begin{quote}
Before the new one: say the French for the handcuffs, then the French for police custody.
\end{quote}

\subsection*{You'll want to know: interroger}
\textbf{interroger} — \textquotedblleft{}to question, to interrogate\textquotedblright{}.

\begin{quote}
Les policiers l'ont interrogé pendant trois heures.
\end{quote}

\begin{quote}
\emph{The police questioned him for three hours.}
\end{quote}

At school, \emph{interroger un élève} — to test a pupil orally.

\begin{cousinweb}
\emph{interroger} is a learned borrowing from Latin \emph{interrogāre} \textquotedblleft{}to ask, to question\textquotedblright{} (EtymDB), built on \emph{rogāre} \textquotedblleft{}to ask\textquotedblright{}.
\end{cousinweb}

\begin{grammarlens}[title={What you've built}]
One more word for telling, explaining and describing at B1.
\end{grammarlens}

\subsection*{Your turn}
\begin{itemize}
\raggedright
  \item \emph{Say it:} \emph{interroger}
  \item \emph{Say it:} \emph{interroger}, once more
  \item \emph{Say it:} a sentence of your own with \emph{interroger}
  \item \emph{Recall:} say the French for the handcuffs, then the French for police custody, then say \emph{interroger} again
\end{itemize}

\begin{tcolorbox}[breakable,colback=teal!4,colframe=teal!35!black,title={Before you move on}]
What does interroger mean? (\textquotedblleft{}To question, to interrogate\textquotedblright{}.) Say it once more.
\end{tcolorbox}
